\section{Conclusion}
One central EEG channel was enough to detect N3 sleep in seven cohorts. Ranking performance (ROC area at least 0.96) was high in every cohort, F1 was lower in the two Dreem cohorts, and accuracy matched DeepSleepNet and exceeded three other published architectures trained under the same protocol, with only a fraction of their parameters. When the detector disagreed with external references, it mostly did so on epochs on which the experts themselves were divided: 51--61\% of the Dreem false positives had at least one expert voting N3. For every detected epoch, transparent factor scores and an additive rating described slow-wave strength, decreased across the night and followed the fraction of experts who scored the epoch as N3, and a reliability score showed which epochs need review. In exploratory post hoc analyses, slow-wave morphology followed the night-time decline more closely than the detector output, and in CAP the scores decreased with age. The three-level grade came close to, but did not reach, its prespecified agreement with expert stage 3/4 scoring. Together, these results show that deep sleep can be detected with high ranking accuracy, described transparently and checked for reliability from a single channel, giving an open and reproducible basis for long-term studies of slow-wave sleep and brain health. The next steps are expert annotation of slow-wave strength within N3 and evaluation on wearable montages. % sensitivity_rq1_false_positive_structure 0.510, 0.609
